\documentclass[oneside,a4paper,14pt]{extarticle} %размер шрифта 14
\usepackage[T1,T2A]{fontenc}
\usepackage[a4paper,% размер страницы
letterpaper,top=20mm,bottom=20mm,
left=15mm,right=15mm]{geometry} %размеры пустых полей
% Можно так \usepackage[left=2.5cm,right=1.5cm,top=2cm,bottom=2cm]{geometry} 
\usepackage[utf8]{inputenc} %кодировка текста
\usepackage[russian]{babel} %поддержка русского языка
\usepackage{textcomp} %текстовые символы
\usepackage{indentfirst} %корректировка отступов
\usepackage{graphicx} %работа с изображениям
\usepackage{mwe} % for blindtext and example-image-a in example
\usepackage{wrapfig}
\usepackage{caption}
\captionsetup[figure]{
name=Рисунок,
labelsep=endash,
justification=centering,
singlelinecheck=false,
}
\usepackage{amsmath}  % для формул и символов
\usepackage{amsfonts}
\usepackage{amsthm}
\usepackage{graphicx}
\usepackage[all]{xy}
\usepackage[breaklinks]{hyperref}
%%размеркегляузаголоковразделов 
\usepackage{titlesec} 
\titleformat{\section}
{\normalsize\bfseries} 
{\thesection} {1em}{} 
\titleformat{\subsection}
{\normalsize\bfseries}
{\thesubsection} {1em}{} 
\titleformat{\subsubsection}
{\normalsize\bfseries}
{\thesubsection} {1em}{}

\renewcommand\baselinestretch{1.33}\normalsize % межстрочный интервал
\setlength{\parindent}{1.25cm}
\usepackage{indentfirst}
\usepackage{float}
\usepackage{listings}
\usepackage[russian]{babel}
\usepackage{listings}
\lstset{
    numberstyle=\tiny,
    basicstyle=\ttfamily\footnotesize,
    breakatwhitespace=false,         
    breaklines=true,                 
    captionpos=b,                    
    keepspaces=true,                 
    numbers=left,                    
    numbersep=5pt,                  
    showspaces=false,                
    showstringspaces=false,
    showtabs=false,                  
    tabsize=2,
    lineskip=1pt,
    extendedchars=true,
    inputencoding=utf8,
    literate={а}{{\selectfont\char224}}1
            {б}{{\selectfont\char225}}1
            {в}{{\selectfont\char226}}1
            {г}{{\selectfont\char227}}1
            {д}{{\selectfont\char228}}1
            {е}{{\selectfont\char229}}1
            {ё}{{\"e}}1
            {ж}{{\selectfont\char230}}1
            {з}{{\selectfont\char231}}1
            {и}{{\selectfont\char232}}1
            {й}{{\selectfont\char233}}1
            {к}{{\selectfont\char234}}1
            {л}{{\selectfont\char235}}1
            {м}{{\selectfont\char236}}1
            {н}{{\selectfont\char237}}1
            {о}{{\selectfont\char238}}1
            {п}{{\selectfont\char239}}1
            {р}{{\selectfont\char240}}1
            {с}{{\selectfont\char241}}1
            {т}{{\selectfont\char242}}1
            {у}{{\selectfont\char243}}1
            {ф}{{\selectfont\char244}}1
            {х}{{\selectfont\char245}}1
            {ц}{{\selectfont\char246}}1
            {ч}{{\selectfont\char247}}1
            {ш}{{\selectfont\char248}}1
            {щ}{{\selectfont\char249}}1
            {ъ}{{\selectfont\char250}}1
            {ы}{{\selectfont\char251}}1
            {ь}{{\selectfont\char252}}1
            {э}{{\selectfont\char253}}1
            {ю}{{\selectfont\char254}}1
            {я}{{\selectfont\char255}}1
            {А}{{\selectfont\char192}}1
            {Б}{{\selectfont\char193}}1
            {В}{{\selectfont\char194}}1
            {Г}{{\selectfont\char195}}1
            {Д}{{\selectfont\char196}}1
            {Е}{{\selectfont\char197}}1
            {Ё}{{\"E}}1
            {Ж}{{\selectfont\char198}}1
            {З}{{\selectfont\char199}}1
            {И}{{\selectfont\char200}}1
            {Й}{{\selectfont\char201}}1
            {К}{{\selectfont\char202}}1
            {Л}{{\selectfont\char203}}1
            {М}{{\selectfont\char204}}1
            {Н}{{\selectfont\char205}}1
            {О}{{\selectfont\char206}}1
            {П}{{\selectfont\char207}}1
            {Р}{{\selectfont\char208}}1
            {С}{{\selectfont\char209}}1
            {Т}{{\selectfont\char210}}1
            {У}{{\selectfont\char211}}1
            {Ф}{{\selectfont\char212}}1
            {Х}{{\selectfont\char213}}1
            {Ц}{{\selectfont\char214}}1
            {Ч}{{\selectfont\char215}}1
            {Ш}{{\selectfont\char216}}1
            {Щ}{{\selectfont\char217}}1
            {Ъ}{{\selectfont\char218}}1
            {Ы}{{\selectfont\char219}}1
            {Ь}{{\selectfont\char220}}1
            {Э}{{\selectfont\char221}}1
            {Ю}{{\selectfont\char222}}1
            {Я}{{\selectfont\char223}}1
}




% ===== добавлено для схем, листингов и скриншотов =====
\usepackage{setspace}   % \singlespacing в таблицах
\usepackage{xcolor}
\usepackage{tikz}
% Общие стили блок-схем (ГОСТ 19.701 / ЕСПД): оформление как в образце отчёта
\usetikzlibrary{arrows.meta,shapes.geometric,shapes.misc,calc,positioning}
\newcommand{\flowfont}{\sffamily\footnotesize}
\tikzset{
  flowbase/.style={draw, line width=1.1pt, font=\flowfont, align=center, inner sep=3pt},
  term/.style={flowbase, rounded rectangle, rounded rectangle arc length=180,
               minimum width=4.4cm, minimum height=1.0cm},
  proc/.style={flowbase, rectangle, minimum width=4.4cm, minimum height=1.0cm, text width=4.0cm},
  io/.style={flowbase, trapezium, trapezium left angle=68, trapezium right angle=112,
             trapezium stretches=true, minimum width=4.6cm, minimum height=1.0cm, text width=3.3cm},
  dec/.style={flowbase, diamond, aspect=3.4, minimum width=4.4cm, minimum height=1.3cm, inner sep=0pt},
  loopS/.style={flowbase, chamfered rectangle, chamfered rectangle corners={north west,north east},
                chamfered rectangle xsep=6mm, chamfered rectangle ysep=6mm,
                minimum width=4.4cm, minimum height=1.1cm, text width=3.8cm},
  loopE/.style={flowbase, chamfered rectangle, chamfered rectangle corners={south west,south east},
                chamfered rectangle xsep=6mm, chamfered rectangle ysep=6mm,
                minimum width=4.4cm, minimum height=1.1cm, text width=3.8cm},
  flow/.style={line width=1.1pt, -{Stealth[length=2.6mm,width=2.2mm]}},
  line/.style={line width=1.1pt},
}
% Подпись «Да»/«Нет» у стрелки: \lab{узел-источник.якорь}{смещение}{текст}
\newcommand{\lab}[3]{\node[font=\sffamily\scriptsize, inner sep=1pt] at ($(#1)+(#2)$) {#3};}
% Комментарий с пунктирной линией и скобкой: \cmt[x, см]{якорь}{имя}{текст}
\newcommand{\cmt}[4][5.2]{%
  \path let \p1=(#2) in node[draw=none, font=\sffamily\scriptsize, anchor=west, align=left,
        text width=3.5cm, inner xsep=7pt] (#3) at (#1cm,\y1) {#4};
  \draw[dashed, line width=0.9pt] (#2) -- (#3.west);
  \draw[line width=1.1pt] ([xshift=5pt]#3.north west) -- (#3.north west)
                          -- (#3.south west) -- ([xshift=5pt]#3.south west);}
       % стили блок-схем (ГОСТ)
% символы в листингах, которых нет в таблице literate выше (добавляются к ней, а не заменяют)
\lstset{literate=*{«}{{\guillemotleft}}1 {»}{{\guillemotright}}1 {—}{{\textemdash}}1 {→}{{\ensuremath{\rightarrow}}}1}

% рисунок-блок-схема из папки fig/: \flowfig[масштаб]{файл}{подпись}
\newcommand{\flowfig}[3][1]{%
  \begin{figure}[H]\centering
  \scalebox{#1}{\input{fig/#2}}
  \caption{#3}\end{figure}}

% листинг из папки code/: \code{язык}{файл}{заголовок}
\newcommand{\code}[3]{%
  \par\noindent\textbf{#3}\par\vspace{2pt}
  \lstinputlisting[language={#1}]{code/#2}}

% скриншот из папки screens/; если файла нет, выводится рамка-заглушка
\newcommand{\screenshot}[3]{% #1 имя файла без .png, #2 подпись, #3 что должно быть на снимке
  \begin{figure}[H]\centering
  \IfFileExists{screens/#1.png}{%
    \includegraphics[width=0.9\textwidth,height=0.42\textheight,keepaspectratio]{screens/#1.png}}{%
    \fbox{\parbox{0.8\textwidth}{\centering\vspace{1.2cm}\colorbox{yellow}{[ скриншот: #3 ]}\\[2pt]{\footnotesize файл screens/#1.png}\vspace{1.2cm}}}}
  \caption{#2}\end{figure}}

\newcommand{\repourl}{https://github.com/lokiPrt22/REPO}   % <- замените на ссылку на свой репозиторий
% ===== конец добавленного =====

\begin{document}
    \newpage\thispagestyle{empty}
    \begin{center}
        МИНИСТЕРСТВО НАУКИ И ВЫСШЕГО ОБРАЗОВАНИЯ\\
            РОССИЙСКОЙ ФЕДЕРАЦИИ
            ФЕДЕРАЛЬНОЕ ГОСУДАРСТВЕННОЕ БЮДЖЕТНОЕ\\
            ОБРАЗОВАТЕЛЬНОЕ
            УЧРЕЖДЕНИЕ ВЫСШЕГО ОБРАЗОВАНИЯ\\
            «ВЯТСКИЙ ГОСУДАРСТВЕННЫЙ УНИВЕРСИТЕТ»\\
            Институт математики и информационных систем\\
            Факультет автоматики и вычислительной техники\\
            Кафедра электронных вычислительных машин
    \end{center}
    \vspace{20mm}
    
    \begin{center}
        Отчёт по лабораторной работе №4\\
        по дисциплине\\
        <<Программирование>>\\
        <<Динамические структуры данных>>\\
        Часть 2\\
    \end{center}
    \vspace{48mm}
    
    Выполнил студент гр. ИВТб-1304-06-00 \hspace{2mm} \rule[-0,5mm]{17mm}{0.15mm}\,/Байбородин Н.А./
    
    
    Проверил старший преподаватель кафедры ЭВМ \hspace{2mm} \rule[-0,5mm]{17mm}{0.15mm}\,/Крутиков А.К./
    
    \vfill
    \begin{center}
        Киров\\
        2026
    \end{center}

    \newpage\thispagestyle{empty}

\section{Цель}
Цель работы: научиться создавать простые графические интерфейсы с использованием библиотеки tkinter, организовывать обработку событий и интегрировать ранее разработанные программные модули в единое приложение. Закрепить понимание работы выбранной структуры данных через её визуальное представление.

\section{Задание}
\begin{itemize}
\item На основе реализации структуры данных из лабораторной работы 4.1 необходимо разработать оконное приложение с графическим интерфейсом на Python.
\item Все операции со структурой данных должны выполняться исключительно через вызовы функций отдельно оформленных библиотечных модулей.
\item Необходимо предусмотреть реализацию как минимум двух модулей: один на C++, второй на Python.
\item Для получения дополнительных баллов необходима реализация двух вариантов C++ модуля: с использованием структур на динамической памяти и с использованием контейнеров и адаптеров STL.
\item Интерфейс должен адекватно реагировать на попытку выполнить некорректную операцию (например, удаление из пустой очереди), выводя понятное сообщение, без аварийного завершения программы.
\item Визуализация должна обновляться после каждой успешной операции.
\end{itemize}
Вариант №1: структура данных «очередь» (как в лабораторной работе 4.1).

% ======================= Теория =======================
\clearpage
\section{Структура FIFO}
Очередь представляет собой линейную структуру данных, организованную по принципу «первым пришёл~--- первым ушёл» (First In, First Out). Это аналогично обычной очереди в магазине: первый покупатель обслуживается первым, а новый становится в конец. Основные операции: Enqueue (добавление в конец), Dequeue (удаление из начала), Peek (просмотр первого элемента без удаления) и Clear (очистка).

\section{Выделение памяти}
В C++ модуле на динамической памяти узлы очереди выделяются оператором \texttt{new} (в варианте \texttt{nothrow}, чтобы нехватка памяти не приводила к исключению) и освобождаются оператором \texttt{delete}. При добавлении элемента создаётся новый узел, при удалении~--- память узла возвращается системе, поэтому память расходуется ровно под текущее количество элементов. В варианте с STL выделением памяти управляет контейнер \texttt{std::deque}, лежащий в основе адаптера \texttt{std::queue}.

\section{Хранение элементов}
Очередь в модуле на динамической памяти реализована на основе односвязного списка. Каждый узел содержит два поля:
\begin{itemize}
\item \texttt{data}~--- хранимое значение (тип \texttt{int}, занимает 4 байта);
\item \texttt{next}~--- указатель на следующий узел (на 64-битной системе занимает 8 байт).
\end{itemize}
Для управления очередью используются два указателя и счётчик:
\begin{itemize}
\item \texttt{head}~--- указывает на первый узел (откуда удаляем);
\item \texttt{tail}~--- указывает на последний узел (куда добавляем);
\item \texttt{count}~--- число элементов в очереди.
\end{itemize}

\begin{figure}[H]\centering
\scalebox{0.82}{\input{fig/queue_struct}}
\caption{Схематичное изображение структуры очереди}
\end{figure}
Общий размер одного узла составляет 16 байт с учётом выравнивания памяти компилятором.

\section{Сложность операций}
\begin{itemize}
\item Добавление элемента~--- $O(1)$;
\item удаление элемента~--- $O(1)$;
\item просмотр первого элемента~--- $O(1)$;
\item очистка и получение всех элементов для отрисовки~--- $O(n)$.
\end{itemize}
Такая реализация обеспечивает высокую скорость работы за счёт прямого доступа к началу и концу очереди через указатели \texttt{head} и \texttt{tail}.

% ======================= Структура приложения =======================
\clearpage
\section{Структура приложения}
Приложение состоит из графического интерфейса и трёх взаимозаменяемых реализаций очереди с одинаковым набором операций: \texttt{enqueue}, \texttt{dequeue}, \texttt{peek}, \texttt{size}, \texttt{clear}, \texttt{items}, \texttt{close}. Интерфейс не знает, какая именно реализация используется: объект очереди создаётся фабрикой \texttt{create\_queue}, а все операции выполняются только через методы этого объекта, то есть через функции библиотечных модулей.

\begin{figure}[H]\centering
\scalebox{0.72}{\input{fig/app_struct}}
\caption{Структура приложения}
\end{figure}

\begin{table}[H]
\caption{Файлы проекта}
\centering\small\singlespacing
\begin{tabular}{|p{3.6cm}|p{11.4cm}|}\hline
\textbf{Файл} & \textbf{Назначение}\\\hline
\texttt{queue\_dyn.cpp} & C++ модуль №1: очередь на динамической памяти (односвязный список, \texttt{head}, \texttt{tail}, \texttt{count}), перенос реализации из лабораторной 4.1\\\hline
\texttt{queue\_stl.cpp} & C++ модуль №2: очередь на адаптере STL \texttt{std::queue<int>}\\\hline
\texttt{queue\_py.py} & Python модуль: очередь на односвязном списке\\\hline
\texttt{queue\_cpp.py} & Обёртка \texttt{ctypes} над C++ библиотеками\\\hline
\texttt{backends.py} & Фабрика реализаций (\texttt{create\_queue}, \texttt{available\_backends})\\\hline
\texttt{app.py} & Графический интерфейс на tkinter\\\hline
\texttt{test\_backends.py} & Одинаковые тесты для всех реализаций\\\hline
\texttt{build.bat}, \texttt{Makefile} & Сборка C++ библиотек\\\hline
\end{tabular}
\end{table}

\begin{table}[H]
\caption{Функции C-интерфейса, экспортируемые обеими C++ библиотеками}
\centering\small\singlespacing
\begin{tabular}{|p{5.2cm}|p{9.8cm}|}\hline
\textbf{Функция} & \textbf{Описание}\\\hline
\texttt{void* q\_create()} & создаёт очередь, возвращает указатель на неё\\\hline
\texttt{void q\_destroy(void*)} & освобождает всю память очереди\\\hline
\texttt{int q\_enqueue(void*, int)} & добавляет элемент в конец: 1~--- успех, 0~--- нехватка памяти\\\hline
\texttt{int q\_dequeue(void*, int*)} & удаляет первый элемент: 1~--- успех, 0~--- очередь пуста\\\hline
\texttt{int q\_peek(void*, int*)} & читает первый элемент: 1~--- успех, 0~--- очередь пуста\\\hline
\texttt{int q\_size(void*)} & возвращает число элементов\\\hline
\texttt{void q\_clear(void*)} & удаляет все элементы\\\hline
\texttt{int q\_to\_array(void*, int*, int)} & копирует элементы от начала к концу в буфер, возвращает их число\\\hline
\end{tabular}
\end{table}
Различие способов сообщения об ошибке скрыто в обёртке \texttt{queue\_cpp.py}: код возврата 0 C++ функций преобразуется в исключение \texttt{QueueEmptyError}, а Python-реализация выбрасывает то же исключение сама. Поэтому интерфейс обрабатывает ошибки одинаково для любой реализации. Модуль на STL отличается только внутренним устройством: \texttt{q\_enqueue} вызывает \texttt{push}, \texttt{q\_dequeue}~--- \texttt{front} и \texttt{pop}, \texttt{q\_clear}~--- обмен с пустой очередью, а так как \texttt{std::queue} не допускает обхода, \texttt{q\_to\_array} выполняется на копии адаптера.

% ======================= Схемы =======================
\clearpage
\section{Схемы алгоритма}
Сначала приведены схемы функций библиотечного модуля на динамической памяти, затем схемы работы графического интерфейса. Схемы операций Python модуля и модуля на STL по логике совпадают с приведёнными.

\flowfig[0.78]{enqueue}{Схема алгоритма для функции q\_enqueue}
\flowfig[0.78]{dequeue}{Схема алгоритма для функции q\_dequeue}
\flowfig{peek}{Схема алгоритма для функции q\_peek}
\flowfig[0.9]{toarray}{Схема алгоритма для функции q\_to\_array (вывод очереди)}
\flowfig[0.9]{clear}{Схема алгоритма для функции q\_clear}
\flowfig[0.78]{main}{Схема алгоритма главного блока программы}
\flowfig[0.75]{on_enqueue}{Схема алгоритма для обработчика on\_enqueue}
\flowfig{on_dequeue}{Схема алгоритма для обработчика on\_dequeue}
\flowfig{on_peek}{Схема алгоритма для обработчика on\_peek}
\flowfig[0.9]{on_clear}{Схема алгоритма для обработчика on\_clear}
\flowfig[0.72]{redraw}{Схема алгоритма для функции redraw (визуализация очереди)}
\flowfig[0.8]{switch}{Схема алгоритма для функции \_switch\_backend (смена реализации)}

% ======================= Обработка ошибок =======================
\section{Обработка некорректных операций}
\begin{table}[H]
\caption{Реакция программы на некорректные действия}
\centering\small\singlespacing
\begin{tabular}{|p{5.6cm}|p{9.4cm}|}\hline
\textbf{Ситуация} & \textbf{Реакция программы}\\\hline
Dequeue или Peek для пустой очереди & предупреждение «Очередь пуста\ldots{}», состояние очереди не меняется\\\hline
Пустое поле ввода & сообщение «Введите целое число.»\\\hline
Ввод не числа & сообщение «\ldots{} не является целым числом.»\\\hline
Число вне диапазона 32-битного \texttt{int} & сообщение «Число не помещается в 32-битный int.»\\\hline
Clear для пустой очереди & сообщение «Очередь уже пуста.»\\\hline
Clear непустой очереди & запрос подтверждения; при отказе очередь не меняется\\\hline
Не найдена C++ библиотека & реализация не показывается в списке, остаётся Python реализация\\\hline
\end{tabular}
\end{table}
Все сообщения выводятся в окне сообщения и дублируются в журнале и строке состояния, аварийного завершения программы не происходит.

\section{Сборка и запуск}
Для сборки C++ библиотек в Windows нужен компилятор MinGW-w64 (\texttt{g++} в переменной \texttt{PATH}, разрядность совпадает с разрядностью Python). Сборка выполняется командой \texttt{build.bat}, в Linux и macOS~--- командой \texttt{make}. Ключ \texttt{-static} нужен, чтобы библиотеки не зависели от динамических библиотек MinGW: Python при загрузке библиотеки их не находит. Запуск приложения: \texttt{python app.py}, запуск тестов: \texttt{python test\_backends.py}.

% ======================= Код =======================
\clearpage
\section{Код интерфейса программы}
\lstinputlisting[language=Python]{code/app.py}

\clearpage
\section{Код фабрики реализаций и обёртки ctypes}
\code{Python}{backends.py}{backends.py}
\code{Python}{queue_cpp.py}{queue\_cpp.py}

\section{Код модуля на Python}
\lstinputlisting[language=Python]{code/queue_py.py}

\section{Код модуля на C++ (динамическая структура данных)}
\lstinputlisting[language=C++]{code/queue_dyn.cpp}

\section{Код модуля на C++ (контейнеры и адаптеры STL)}
\lstinputlisting[language=C++]{code/queue_stl.cpp}

\section{Код тестов и сборки}
\code{Python}{test_backends.py}{test\_backends.py}
\code{}{build.bat}{build.bat}
\code{make}{Makefile}{Makefile}

% ======================= Тестирование =======================
\clearpage
\section{Тестирование}
Файл \texttt{test\_backends.py} выполняет один и тот же набор проверок для каждой доступной реализации: пустая очередь (Dequeue и Peek вызывают ошибку), порядок FIFO, Peek не удаляет элемент, опустошение и повторное использование очереди (сброс указателя \texttt{tail}), очистка, добавление 1000 элементов. Результат выполнения приведён на рисунке ниже.

\screenshot{tests}{Результат выполнения тестов}{вывод команды python test\_backends.py, строка OK}

% ======================= Экранные формы =======================
\section{Экранные формы}
\begin{figure}[H]\centering
\includegraphics[width=0.9\textwidth]{screens/select.png}
\caption{Главное окно приложения, выбор реализации очереди}
\end{figure}
\screenshot{filled}{Очередь с несколькими элементами}{окно с 3--4 элементами, видны Front и Rear}
\screenshot{peek}{Операция Peek}{первый элемент подсвечен зелёным}
\screenshot{dequeue}{Операция Dequeue}{очередь после удаления, запись в журнале}
\screenshot{errempty}{Сообщение об ошибке при удалении из пустой очереди}{предупреждение «Очередь пуста»}
\screenshot{errinput}{Сообщение об ошибке при некорректном вводе}{сообщение при вводе текста вместо числа}
\screenshot{switch}{Переключение реализации очереди}{выбрана C++ (STL std::queue), элементы сохранились}

% ======================= Github =======================
\section{Github}
\screenshot{github}{Снимок экрана репозитория}{страница репозитория на GitHub со списком файлов}
Ссылка~--- \url{\repourl}

% ======================= Вывод =======================
\clearpage
\section{Вывод}
В ходе выполнения лабораторной работы было разработано оконное приложение с графическим интерфейсом на Python с использованием библиотеки tkinter, визуализирующее структуру данных «очередь» из лабораторной работы 4.1. Реализована интеграция трёх библиотечных модулей: двух модулей на C++ (на динамической памяти и на адаптере STL \texttt{std::queue}) и модуля на Python. Все операции над очередью интерфейс выполняет только через функции этих модулей, а реализации можно менять прямо во время работы программы.

Графический интерфейс обеспечивает выполнение операций Enqueue, Dequeue, Peek и Clear и делает взаимодействие со структурой данных наглядным: очередь отображается на Canvas в виде цепочки элементов с указателями Front и Rear и перерисовывается после каждой успешной операции. Приложение корректно обрабатывает некорректные действия (удаление из пустой очереди, ввод не числа, переполнение типа \texttt{int}), выводя понятные сообщения и не завершаясь аварийно.

Таким образом, поставленная цель достигнута: закреплены практические навыки разработки GUI-приложений, обработки событий, организации взаимодействия между компонентами программы и интеграции разноязыковых модулей в единое приложение.

\end{document}
